% Canny 两步后处理的教学例；数值与正文算例一致。左：梯度方向为水平时的非极大值抑制；
% 右：沿一条轮廓的八个候选点与双阈值（T_high=40，T_low=20）。宽约 98mm。
\begin{tikzpicture}[x=1mm,y=1mm,font=\figurefont\small,
  cell/.style={draw=FigureStroke,line width=.5pt,minimum size=8mm,inner sep=0pt},
  cmp/.style={cell,fill=FigureBlueFill,draw=FigureBlue,line width=.8pt},
  ctr/.style={cell,fill=FigureYellowFill,draw=FigureYellow,line width=.9pt},
  lab/.style={font=\figurefont\footnotesize,text=FigureMuted}]
  % (a) NMS
  \foreach \r/\row in {0/{30,45,28},1/{80,120,100},2/{35,50,40}}{
    \foreach \v [count=\c from 0] in \row{
      \node[cell] at (\c*8,-\r*8) {\v};
    }
  }
  \node[cmp] at (0,-8) {80};
  \node[ctr] at (8,-8) {120};
  \node[cmp] at (16,-8) {100};
  \draw[-{Stealth[length=2mm]},FigureStroke,line width=.9pt] (-2,-21)--(18,-21);
  \node[lab,anchor=north] at (8,-22.5) {梯度方向};
  \node at (8,8) {(a) 非极大值抑制};
  % (b) hysteresis
  \begin{scope}[xshift=33mm,yshift=-26mm,y=.42mm]
    \foreach \v [count=\k from 0] in {15,42,38,25,18,33,60,12}{
      \pgfmathsetmacro{\keep}{(\k==1||\k==2||\k==3||\k==5||\k==6)?1:0}
      \ifdim\keep pt>0pt
        \fill[FigureBlueFill,draw=FigureBlue,line width=.7pt] (\k*7.4-3,0) rectangle (\k*7.4+3,\v);
      \else
        \draw[FigureMuted,line width=.6pt] (\k*7.4-3,0) rectangle (\k*7.4+3,\v);
      \fi
      \node[font=\figurefont\footnotesize,anchor=south] at (\k*7.4,\v) {\v};
      \node[font=\figurefont\footnotesize,anchor=north] at (\k*7.4,-1) {$\k$};
    }
    \draw[FigureStroke,line width=.6pt] (-5,0)--(57,0);
    \draw[FigureRed,dashed,line width=.7pt] (-5,40)--(57,40) node[right,font=\figurefont\footnotesize,text=FigureStroke] {$T_{\mathrm{high}}$};
    \draw[FigureYellow,dashed,line width=.7pt] (-5,20)--(57,20) node[right,font=\figurefont\footnotesize,text=FigureStroke] {$T_{\mathrm{low}}$};
    \node at (26,82) {(b) 滞后阈值};
  \end{scope}
\end{tikzpicture}
